\documentclass[10pt]{extarticle}
\usepackage[a4paper,margin=1.5cm]{geometry}
\usepackage[T1]{fontenc}
\usepackage{amsmath,amssymb}
\usepackage{graphicx}
\usepackage[export]{adjustbox}
\usepackage{multicol}
\usepackage{enumitem}
\usepackage{booktabs,array,tabularx}
\setlength{\columnsep}{0.7cm}
\setlength{\parindent}{0pt}
\setlength{\parskip}{3pt}
\pagestyle{empty}
\makeatletter
% figure registry (figdefs.tex is generated by the extract step)
\newcommand{\qfigdef}[4]{\@namedef{qf@file@#1}{#2}\@namedef{qf@w@#1}{#3}\@namedef{qf@wide@#1}{#4}}
\newcommand{\qf@check}[1]{\@ifundefined{qf@file@#1}{\PackageError{score2}{Unknown figure id #1}{}}{}}
\newcommand{\qf@img}[4][]{\includegraphics[width=#2\dimexpr\@nameuse{qf@w@#4}pt\relax,max width=#3,max height=0.45\textheight,#1]{\@nameuse{qf@file@#4}}}
% Full-width questions: finalize wraps every question that contains a wide figure in \qfullbegin...\qfullend,
% so the WHOLE question (stem, figures, options) is set across the page in reading order, then the two columns
% resume. The question's head - its number, stem and first figure (\qfullsplit marks the end of the head) -
% is kept on one page, so a number is never left at a page foot with its figure overleaf; the rest may break
% across pages normally, so no big gaps are left. Without a \qfullsplit, the whole question is kept together
% if it is under 45% of a page. Never write these yourself.
\newif\ifqf@full
\newif\ifqf@boxed
\newsavebox{\qf@box}
\newcommand{\qf@open}{\global\qf@fulltrue\par\global\qf@boxedtrue
  \setbox\qf@box\vbox\bgroup\hsize\textwidth\linewidth\textwidth}
\newcommand{\qf@emit}[1]{\par\egroup\global\qf@boxedfalse
  \ifdim\dimexpr\ht\qf@box+\dp\qf@box\relax<#1\textheight\noindent\usebox{\qf@box}\par\else\unvbox\qf@box\fi}
\newcommand{\qfullsplit}{\ifqf@boxed\qf@emit{0.85}\fi}
\newcommand{\qf@close}{\ifqf@boxed\qf@emit{0.45}\fi\global\qf@fullfalse}
\newcommand{\qfullbegin}{\end{multicols}\qf@open}
\newcommand{\qfullend}{\qf@close\begin{multicols}{2}}
\newcommand{\qfullstart}{\qf@open}
\newcommand{\qfulldone}{\qf@close}
% \qfig[scale]{id}: block figure, centred, directly where it is written (size set by the tool).
% Fallback if a wide figure is not inside a full-width question: break out of the columns in place - no floats.
\newcommand{\qfig}[2][1]{\qf@check{#2}%
  \ifqf@full\qf@blk{#1}{#2}\else
  \ifnum\@nameuse{qf@wide@#2}=1 \expandafter\@firstoftwo\else\expandafter\@secondoftwo\fi
  {\end{multicols}\begin{center}\qf@img{#1}{\textwidth}{#2}\end{center}\begin{multicols}{2}}%
  {\qf@blk{#1}{#2}}\fi}
\newcommand{\qf@blk}[2]{\par\smallskip{\centering\qf@img{#1}{\linewidth}{#2}\par}\smallskip}
% \qfigi[scale]{id}: inline figure (inside option lists / tables / minipages); never breaks columns.
\newcommand{\qfigi}[2][1]{\qf@check{#2}\qf@img[valign=c]{#1}{\linewidth}{#2}}
\makeatother
% structure macros - lint reads these to check completeness and figure placement
\newcommand{\papersec}[1]{\par\bigskip{\centering\Large\bfseries PAPER-#1\par}\medskip}
\newcommand{\partsec}[2]{\par\medskip\noindent{\large\bfseries PART-#1: #2}\par\smallskip}
\newcommand{\sect}[1]{\par\smallskip\noindent\textbf{#1}\par}
\newcommand{\qn}[1]{\par\smallskip\noindent\textbf{#1.}~}
\newcommand{\paraq}[2]{\par\smallskip\noindent\textbf{Paragraph for Questions #1 and #2}\par}
\newenvironment{qoptions}{\begin{itemize}[leftmargin=1.8em,itemsep=1pt,topsep=2pt,label={}]}{\end{itemize}}
\newcommand{\opt}[2]{\item[(#1)] #2}
% figure registry (generated by score2_tools.py finalize)
\qfigdef{p1_phy_s1_q01_1}{Score Advanced - Test 35 (15 May 2024)-img/p1_phy_s1_q01_1.png}{101.3}{0}
\qfigdef{p1_phy_s1_q02_1}{Score Advanced - Test 35 (15 May 2024)-img/p1_phy_s1_q02_1.png}{101.8}{0}
\qfigdef{p1_phy_s1_q03_1}{Score Advanced - Test 35 (15 May 2024)-img/p1_phy_s1_q03_1.png}{152.6}{0}
\qfigdef{p1_phy_s1_q05_1}{Score Advanced - Test 35 (15 May 2024)-img/p1_phy_s1_q05_1.png}{145.2}{0}
\qfigdef{p1_phy_s1_q06_1}{Score Advanced - Test 35 (15 May 2024)-img/p1_phy_s1_q06_1.png}{74.2}{0}
\qfigdef{p1_phy_s2_q07_1}{Score Advanced - Test 35 (15 May 2024)-img/p1_phy_s2_q07_1.png}{108.3}{0}
\qfigdef{p1_phy_s2_q09_1}{Score Advanced - Test 35 (15 May 2024)-img/p1_phy_s2_q09_1.png}{138.7}{0}
\qfigdef{p1_phy_s2_q10_1}{Score Advanced - Test 35 (15 May 2024)-img/p1_phy_s2_q10_1.png}{119.2}{0}
\qfigdef{p1_phy_s3_q01_1}{Score Advanced - Test 35 (15 May 2024)-img/p1_phy_s3_q01_1.png}{85.8}{0}
\qfigdef{p1_phy_s3_q05_1}{Score Advanced - Test 35 (15 May 2024)-img/p1_phy_s3_q05_1.png}{72.8}{0}
\qfigdef{p1_phy_s3_q06_1}{Score Advanced - Test 35 (15 May 2024)-img/p1_phy_s3_q06_1.png}{103.1}{0}
\qfigdef{p1_chem_s1_q05_A}{Score Advanced - Test 35 (15 May 2024)-img/p1_chem_s1_q05_A.png}{245.2}{0}
\qfigdef{p1_chem_s1_q05_B}{Score Advanced - Test 35 (15 May 2024)-img/p1_chem_s1_q05_B.png}{266.5}{1}
\qfigdef{p1_chem_s1_q05_C}{Score Advanced - Test 35 (15 May 2024)-img/p1_chem_s1_q05_C.png}{234.4}{0}
\qfigdef{p1_chem_s1_q05_D}{Score Advanced - Test 35 (15 May 2024)-img/p1_chem_s1_q05_D.png}{260.1}{1}
\qfigdef{p1_chem_s3_q06_1}{Score Advanced - Test 35 (15 May 2024)-img/p1_chem_s3_q06_1.png}{143.2}{0}
\qfigdef{p1_chem_s3_q07_1}{Score Advanced - Test 35 (15 May 2024)-img/p1_chem_s3_q07_1.png}{143.0}{0}
\qfigdef{p1_chem_s3_q08_1}{Score Advanced - Test 35 (15 May 2024)-img/p1_chem_s3_q08_1.png}{105.5}{0}
\qfigdef{p1_chem_s3_q08_2}{Score Advanced - Test 35 (15 May 2024)-img/p1_chem_s3_q08_2.png}{237.7}{0}
\qfigdef{p2_phy_s1_q02_A}{Score Advanced - Test 35 (15 May 2024)-img/p2_phy_s1_q02_A.png}{76.1}{0}
\qfigdef{p2_phy_s1_q02_B}{Score Advanced - Test 35 (15 May 2024)-img/p2_phy_s1_q02_B.png}{76.2}{0}
\qfigdef{p2_phy_s1_q02_C}{Score Advanced - Test 35 (15 May 2024)-img/p2_phy_s1_q02_C.png}{72.7}{0}
\qfigdef{p2_phy_s1_q02_D}{Score Advanced - Test 35 (15 May 2024)-img/p2_phy_s1_q02_D.png}{73.7}{0}
\qfigdef{p2_phy_s1_q04_1}{Score Advanced - Test 35 (15 May 2024)-img/p2_phy_s1_q04_1.png}{271.3}{1}
\qfigdef{p2_phy_s2_q05_1}{Score Advanced - Test 35 (15 May 2024)-img/p2_phy_s2_q05_1.png}{119.5}{0}
\qfigdef{p2_phy_s2_q06_1}{Score Advanced - Test 35 (15 May 2024)-img/p2_phy_s2_q06_1.png}{97.2}{0}
\qfigdef{p2_phy_s2_q07_1}{Score Advanced - Test 35 (15 May 2024)-img/p2_phy_s2_q07_1.png}{83.3}{0}
\qfigdef{p2_phy_s2_q08_1}{Score Advanced - Test 35 (15 May 2024)-img/p2_phy_s2_q08_1.png}{101.1}{0}
\qfigdef{p2_phy_s3_q01_1}{Score Advanced - Test 35 (15 May 2024)-img/p2_phy_s3_q01_1.png}{103.8}{0}
\qfigdef{p2_phy_s3_q03_1}{Score Advanced - Test 35 (15 May 2024)-img/p2_phy_s3_q03_1.png}{245.2}{0}
\qfigdef{p2_phy_s3_q06_1}{Score Advanced - Test 35 (15 May 2024)-img/p2_phy_s3_q06_1.png}{143.3}{0}
\qfigdef{p2_phy_s3_q08_1}{Score Advanced - Test 35 (15 May 2024)-img/p2_phy_s3_q08_1.png}{187.6}{0}
\qfigdef{p2_chem_s1_q02_1}{Score Advanced - Test 35 (15 May 2024)-img/p2_chem_s1_q02_1.png}{121.3}{0}
\qfigdef{p2_chem_s1_q04_1}{Score Advanced - Test 35 (15 May 2024)-img/p2_chem_s1_q04_1.png}{290.7}{1}
\qfigdef{p2_chem_s2_q08_1}{Score Advanced - Test 35 (15 May 2024)-img/p2_chem_s2_q08_1.png}{98.0}{0}
\qfigdef{p2_chem_s2_q08_2}{Score Advanced - Test 35 (15 May 2024)-img/p2_chem_s2_q08_2.png}{96.5}{0}
\qfigdef{p2_chem_s2_q08_3}{Score Advanced - Test 35 (15 May 2024)-img/p2_chem_s2_q08_3.png}{113.7}{0}
\qfigdef{p2_chem_s2_q09_1}{Score Advanced - Test 35 (15 May 2024)-img/p2_chem_s2_q09_1.png}{223.1}{0}
\qfigdef{p2_chem_s2_q10_1}{Score Advanced - Test 35 (15 May 2024)-img/p2_chem_s2_q10_1.png}{332.5}{1}
\raggedcolumns
\makeatletter\@beginparpenalty=10000\makeatother
\newcolumntype{L}[1]{>{\raggedright\arraybackslash}p{#1}}
\begin{document}
\begin{center}{\Large\bfseries Score Advanced --- Test 35 (15 May 2024)}\end{center}
\begin{multicols}{2}
\raggedcolumns
\papersec{1}
\partsec{1}{PHYSICS}
\sect{SECTION-I(i)}
\qn{1} Consider the circuit shown with respective specifications of elements marked in the figure. Capacitor-1 is charged such that charge on it is Q$_0$ and it's left plate is positively charged. While capacitor-2 is uncharged. The switch is closed at t = 0.\par
\qfig{p1_phy_s1_q01_1}
\begin{qoptions}
\opt{A}{Frequency of oscillation of charge on left plate of capacitor-1 is $\dfrac1{\pi\sqrt{2LC}}$}
\opt{B}{Frequency of oscillation of charge on left plate of capacitor-1 is $\dfrac1\pi\sqrt{\dfrac2{LC}}$}
\opt{C}{Maximum current through the inductor is $\dfrac{Q_0}{\sqrt{2LC}}$}
\opt{D}{Maximum current through the inductor is $\dfrac{Q_0}{\sqrt{LC}}$}
\end{qoptions}
\qn{2} In the potentiometer circuit shown here, a potentiomete wire AB of length 100 cm and resistance 8 $\Omega$ is used. While measuring potential difference between points X and Y balancing length (AC) is obtained to be equal to 25 cm.\par
\qfig{p1_phy_s1_q02_1}
\begin{qoptions}
\opt{A}{Internal resistance r (in $\Omega$) of the driving cell is 2$\Omega$.}
\opt{B}{Internal resistance r (in $\Omega$) of the driving cell is 4$\Omega$.}
\opt{C}{Potential difference between XY is 2 volt.}
\opt{D}{Potential difference between XY is 4 volt.}
\end{qoptions}
\qn{3} \qfig{p1_phy_s1_q03_1}
A particle of mass m moving with a velocity v$_0$ enters a stationary object of mass M (free to move) as shown in the diagram. All surfaces are frictionless and system shown in the diagram is lying in a horizontal plane.
\begin{qoptions}
\opt{A}{Time taken by particle of mass m to cross the object of mass M is $\dfrac{2\ell+\pi R}{V_0}$}
\opt{B}{Time taken by particle of mass m to cross the object of mass M is log e$\dfrac{\ell+R}{V_0}$}
\opt{C}{Speed of particle (m) w.r.t block (M) will not change}
\opt{D}{Speed of particle (m) w.r.t block (M) will change}
\end{qoptions}
\qn{4} A coil with inductance L and resistance R is connected to an alternating source. What should be the capacity of a capacitor connected in series with the coil such that active power of circuit does not change.
\begin{qoptions}
\opt{A}{C = $\dfrac1{\omega^2}$}
\opt{B}{C = $\dfrac1{\omega^2L}$}
\opt{C}{C = $\dfrac1{2\omega^2L}$}
\opt{D}{C = $\dfrac{\omega^2}L$}
\end{qoptions}
\qn{5} In a YDSE setup, we plot the phase difference ($\phi$) between both waves at point P on the screen against the angular position ($\theta$) of point P on the screen. The graph is as shown below.\par
\qfig{p1_phy_s1_q05_1}
\begin{qoptions}
\opt{A}{The distance S$_1$S$_2$ = 2$\lambda$}
\opt{B}{There are a total of 4 minima on the screen}
\opt{C}{The first maxima above the centre is at $\theta=\dfrac\pi4$}
\opt{D}{At $\theta=\dfrac\pi6$, intensity is maximum}
\end{qoptions}
\qn{6} A uniform rod of mass M and length $\sqrt3$R is intially stationary and is in vertical position inside a frictionless, fixed hollow sphere of radius R. Now bar is relased from rest from this position.\par
\qfig{p1_phy_s1_q06_1}
\begin{qoptions}
\opt{A}{Find velocity of c.o.m when bar becomes horizontal $\sqrt{\dfrac{9R}2}$}
\opt{B}{Angular velocity of bar when it becomes horizontal is $\sqrt{\dfrac{2g}R}$}
\opt{C}{Normal reaction exerted by sphere at one end of rod when rod is horizontal is mg}
\opt{D}{Normal reaction exerted by sphere at one end of rod when rod is horizontal is 2 mg}
\end{qoptions}
\sect{SECTION-I(ii)}
\qn{7} Consider the situation shown in diagram. Both capacitors are having same capacitance C. Capacitors are initially uncharged. Switch K is closed at time = 0 after a long time capacitors are fully charged. List- I represents value of r and List- II represents charge passing through r during complete chargning process. Now match the followings :\par
\qfig{p1_phy_s2_q07_1}
{\footnotesize\setlength{\tabcolsep}{1.5pt}\noindent\begin{tabular}{|c|L{1.6cm}|c|c|}
\hline
\multicolumn{2}{|c|}{List- I} & \multicolumn{2}{c|}{List- II} \\ \hline
(P) & r = 0 & (1) & $\dfrac{CV}6$ \\ \hline
(Q) & r = R & (2) & $\dfrac{CV}5$ \\ \hline
(R) & r = 2R & (3) & $\dfrac{CV}4$ \\ \hline
(S) & r = 3R & (4) & $\dfrac{CV}3$ \\ \hline
\end{tabular}}\par
\begin{qoptions}
\opt{A}{P $\to$ 4;Q $\to$ 2;R $\to$ 3;S $\to$ 1}
\opt{B}{P $\to$ 1;Q $\to$ 4;R $\to$ 3;S $\to$ 2}
\opt{C}{P $\to$ 1;Q $\to$ 2;R $\to$ 3;S $\to$ 4}
\opt{D}{P $\to$ 2;Q $\to$ 4;R $\to$ 1;S $\to$ 3}
\end{qoptions}
\qn{8} Following is given the equation of a stationary wave (all in SI units) y = (0.06) sin (2$\pi$x) cos (5$\pi$t)\par
{\footnotesize\setlength{\tabcolsep}{1.5pt}\noindent\begin{tabular}{|c|L{3.2cm}|c|c|}
\hline
 & \multicolumn{1}{c|}{List$-$ I} & \multicolumn{2}{c|}{\begin{tabular}[c]{@{}c@{}}List$-$ II\\(Only magnitude in SI units)\end{tabular}} \\ \hline
(P) & Amplitude of constituent wave (in m) & (1) & 0.06 \\ \hline
(Q) & Position of node at x =...... (in m) & (2) & 0.5 \\ \hline
(R) & Position of antinode at x =.... (in m) & (3) & 0.25 \\ \hline
(S) & Amplitude at x = 1/12 m (in m) & (4) & 0.03 \\ \hline
\end{tabular}}\par
\begin{qoptions}
\opt{A}{P $\to$ 4;Q $\to$ 2;R $\to$ 4;S $\to$ 3}
\opt{B}{P $\to$ 1;Q $\to$ 4;R $\to$ 3;S $\to$ 2}
\opt{C}{P $\to$ 4;Q $\to$ 2;R $\to$ 3;S $\to$ 4}
\opt{D}{P $\to$ 2;Q $\to$ 4;R $\to$ 1;S $\to$ 3}
\end{qoptions}
\qn{9} Two charged di-electric spheres A (solid) and B (hollow) of same radius R and masses 2m and m respectively are kept on a horizontal rough surface. The two spheres A and B have charges 2Q and $-$Q respectively and uniformly distributed. They are released from rest when their centres are separated by very large distance. Assuming there is no sliding.\par
\qfig{p1_phy_s2_q09_1}
{\footnotesize\setlength{\tabcolsep}{1.5pt}\noindent\begin{tabular}{|c|L{3.6cm}|c|c|}
\hline
 & \multicolumn{1}{c|}{List-I} & \multicolumn{2}{c|}{List-II} \\ \hline
(P) & Kinetic energy of system just before A and B meet & (1) & $\dfrac{838kQ^2}{2345R}$ \\ \hline
(Q) & Kinetic energy of A just before A and B meet & (2) & $\dfrac{kQ^2}{R}$ \\ \hline
(R) & Kinetic energy of B just before A and B meet & (3) & $\dfrac{25kQ^2}{67R}$ \\ \hline
(S) & Rotational kinetic energy of system just before A and B meet & (4) & $\dfrac{42kQ^2}{67R}$ \\ \hline
\end{tabular}}\par
\begin{qoptions}
\opt{A}{P $\to$ 2;Q $\to$ 3;R $\to$ 4;S $\to$ 1}
\opt{B}{P $\to$ 1;Q $\to$ 4;R $\to$ 3;S $\to$ 2}
\opt{C}{P $\to$ 3;Q $\to$ 2;R $\to$ 1;S $\to$ 4}
\opt{D}{P $\to$ 4;Q $\to$ 1;R $\to$ 2;S $\to$ 3}
\end{qoptions}
\qn{10} A very small current carrying square loop (current I) of side `L' is placed in y-z plane with centre at origin of the coordinate system (shown in figure). In List$-$I the coordinate of the points are given \& in List$-$II magnitude of strength of magnetic field is given. Then :-\par
\qfig{p1_phy_s2_q10_1}
{\footnotesize\setlength{\tabcolsep}{1.5pt}\noindent\begin{tabular}{|c|L{3.4cm}|c|c|}
\hline
 & \multicolumn{1}{c|}{List$-$I} & \multicolumn{2}{c|}{List$-$II} \\ \hline
(P) & At point O (0, 0, 0) & (1) & $\dfrac{\mu_0IL^2}{2\pi a^3}$ \\ \hline
(Q) & At point P$_1$ (a, 0, 0) (here a $>>$ L) & (2) & $\dfrac{2\sqrt2\mu_0I}{\pi L}$ \\ \hline
(R) & At point P$_2$ (0, a, 0) (here a $>>$ L) & (3) & $\dfrac{\mu_0\sqrt5IL^2}{16\pi a^3}$ \\ \hline
(S) & At point P$_3$ (a, a, 0) (here a $>>$ L) & (4) & $\dfrac{\mu_0IL^2}{4\pi a^3}$ \\ \hline
\end{tabular}}\par
\begin{qoptions}
\opt{A}{P $\to$ 4;Q $\to$ 2;R $\to$ 3;S $\to$ 1}
\opt{B}{P $\to$ 1;Q $\to$ 1;R $\to$ 2;S $\to$ 3}
\opt{C}{P $\to$ 4;Q $\to$ 1;R $\to$ 2;S $\to$ 3}
\opt{D}{P $\to$ 2;Q $\to$ 1;R $\to$ 4;S $\to$ 3}
\end{qoptions}
\sect{SECTION-II}
\qn{1} A charge Q is uniformally distributed over a square plate of side a. Another point charge q is placed at a height $\dfrac a2$ from the centre of square plate as shown in diagram. Force acting on the point charge q is $\dfrac{nQq}{\in_0a^2}$ . Find n.\par
\qfig{p1_phy_s3_q01_1}
\qn{2} A cell of emf 1.5 V and internal resistance 0.5 $\Omega$ is connected to a nonlinear conductor in which current varies with voltage as V$^2$ = I. Calculate the current drawn from cell (in ampere).
\qn{3} Light falls on a prism of angle cos$^{-1}\left(\dfrac18\right)$. If the minimum value of refractive index of the prism-material, so that light entering one face may not exit through the other face
\qn{4} One mole of an ideal monatomic gas at temperature T$_0$ expands slowly according to the law P/V = constant. If the final temperature is 2T$_0$, heat supplied to the gas is $\dfrac{RT_0}{x_0}$ . Then find value of x?
\qn{5} A hemispherical shell is fixed. The mass distribution of the shell is uniform. At a point on the symmetry axis, which is exactly at a distance of the radius measured from the centre of the hemisphere O, a small bead is released. The speed of the bead is v$_0$ when it reaches the centre O. The speed of the bead be when it reaches the shell is xv$_0$. Find x\ \ (Neglect the effect of gravity.)\par
\qfig{p1_phy_s3_q05_1}
\qn{6} We stick a inextensible adhesive tape on a tabletop. We grab one end of it and pull it back horizontally and uniformly with a velocity of 6 cm/s. What is the velocity (in cm/s) of the middle of the moving part of the tape?\par
\qfig{p1_phy_s3_q06_1}
\qn{7} Two coils, 1 and 2, have a mutual inductance = M and resistance R each. A current flows in coil 1, which varies with time as : i$_1$ = kt$^2$, where k is a constant and `t' is time. The total charge that has flown through coil 2, between t = 0 and t = T is $\left(\dfrac{\lambda}{-\lambda+1}\right)\dfrac{kmT^2}{R}$. Find $\lambda$.
\qn{8} A small bead of mass m and positive charge +q is mounted on a fixed, thin, smooth horizontal plastic ring of radius R. At t = 0, bead is stationary. Ring lies in x-y plane with centre at (0, 0). Magnetic field $\vec B=\dfrac{E_0t}{r}\hat k$ (where r $\neq$ 0) is switched on, where t denotes time, r distance from z-axis and E$_0$ is a positive constant. Magnitude of normal reaction exerted by ring on bead is qE$_0\lambda$. Find $\lambda$ (Gravity is absent).
\raggedcolumns
\papersec{1}
\partsec{2}{CHEMISTRY}
\sect{SECTION-I(i)}
\qn{1} How many of the following reactions will not result in precipitate formation ?\par
(ignore balancing)
\begin{qoptions}
\opt{A}{FeCl$_3$(aq) + Na$_2$SO$_4$(aq) $\to$}
\opt{B}{Hg(NO$_3$)$_2$(aq) + exNaCN(aq) $\to$}
\opt{C}{NiCl$_2$(aq) + ex NH$_3$(aq) $\to$}
\opt{D}{ZnCl$_2$(aq) + H$_2$S(aq) $\xrightarrow{H^+}$}
\end{qoptions}
\qn{2} Which of the following will result in formation of H$_3$BO$_3$ one of the products ?\par
(ignore balancing)
\begin{qoptions}
\opt{A}{Na$_2$[B$_4$O$_5$(OH)$_4$] + dil. H$_2$SO$_4$ $\to$}
\opt{B}{B$_2$H$_6$ + H$_2$O $\to$}
\opt{C}{B$_3$N$_3$H$_6$ + H$_2$O $\to$}
\opt{D}{Na$_2$B$_4$O$_7$ + NH$_4$Cl $\xrightarrow{\text{fusion}}$}
\end{qoptions}
\qn{3} Which of the following statement(s) is(are) correct ?
\begin{qoptions}
\opt{A}{Smoke and cloud both are aerosol with same dispersion medium.}
\opt{B}{Surface tension of lyophilic sols is less than the surface tension of dispersion medium while that of lyophobic sols is more than dispersion medium.}
\opt{C}{Lyophobic sols can be detected by ultra microscope}
\opt{D}{Colloidal sol shows low magnitude of colligative properties, as compared to true solutions.}
\end{qoptions}
\qn{4} If pH of an aqueous solution of 0.1M monoprotic acid HA is found to be 3 at T kelvin then
\begin{qoptions}
\opt{A}{Degree of dissociation of the acid would be 1 \% (approx.)}
\opt{B}{Osmotic pressure of the solution would be 0.011RT}
\opt{C}{Theoretical molar mass of HA is approximately 1.01 times of its observed molar mass}
\opt{D}{K$_\mathrm b$ of A$^-_{\mathrm{(aq.)}}$ is 10$^{-11}$ M (approx.) at 298K}
\end{qoptions}
\qfullbegin
\qn{5} Among the following, the reaction(s) that favor(s) the formation of the products\ \ at 25$^\circ$C is/are
\begin{qoptions}
\opt{A}{\qfigi{p1_chem_s1_q05_A}}
\opt{B}{\qfigi{p1_chem_s1_q05_B}}
\opt{C}{\qfigi{p1_chem_s1_q05_C}}
\opt{D}{\qfigi{p1_chem_s1_q05_D}}
\end{qoptions}
\qfullend
\qn{6} Among the following, the correct statements is/are:
\begin{qoptions}
\opt{A}{The first pK$_\mathrm a$ of malonic acid is lower than the pK$_\mathrm a$ of acetic acid while its second pK$_\mathrm a$ is higher than the pK$_\mathrm a$\ \ of acetic acid.}
\opt{B}{Picric acid shows positive test with NaHCO$_3$.}
\opt{C}{CH(CN)$_3$ is a stronger acid than HCN.}
\opt{D}{Amides are generally stronger acid than water due to resonance in their conjugate base.}
\end{qoptions}
\sect{SECTION-I(ii)}
\qn{7} {\footnotesize\setlength{\tabcolsep}{1.5pt}\noindent\begin{tabular}{|c|L{2.5cm}|c|L{3.4cm}|}
\hline
\multicolumn{2}{|c|}{List-I (Complex)} & \multicolumn{2}{c|}{List-II} \\ \hline
(P) & [Rh(PPh$_3$)$_3$Cl] & (1) & dxy = dyz = dzx (energy) \\ \hline
(Q) & [CoCl$_4$]$^{2-}$ & (2) & High spin \\ \hline
(R) & Ni(CO)$_4$ & (3) & Low spin \\ \hline
(S) & [Fe(H$_2$O)$_5$NO]$^{2+}$ & (4) & Diamagnetic \\ \hline
\end{tabular}}\par
\begin{qoptions}
\opt{A}{P $\to$ 2;Q $\to$ 4;R $\to$ 3;S $\to$ 1}
\opt{B}{P $\to$ 3;Q $\to$ 1;R $\to$ 2;S $\to$ 4}
\opt{C}{P $\to$ 3;Q $\to$ 2;R $\to$ 4;S $\to$ 1}
\opt{D}{P $\to$ 4;Q $\to$ 1;R $\to$ 2;S $\to$ 3}
\end{qoptions}
\qn{8} Match the list\par
{\footnotesize\setlength{\tabcolsep}{1.5pt}\noindent\begin{tabular}{|c|L{3.0cm}|c|L{3.2cm}|}
\hline
 & \multicolumn{1}{c|}{List-I} & & \multicolumn{1}{c|}{List-II} \\ \hline
(P) & Electrolysis of 100 L aqueous solution of CH$_3$COOK by passing 2 F of electricity & (1) & pH = 12.3\newline Anode = Ethane (g) + CO$_2$(g)\newline Cathode = H$_2$(g) \\ \hline
(Q) & Electrolysis of 10 L aqueous solution of HCOOK by passing 1 F of electricity. & (2) & pH = 13.0\newline Anode = H$_2$ (g) + CO$_2$(g)\newline Cathode = H$_2$(g) \\ \hline
(R) & Electrolysis of 10 L aqueous solution of K$_2$SO$_4$ by passing 1 F of electricity. & (3) & pH = 7.0\newline Anode = O$_2$(g)\newline Cathode = H$_2$(g) \\ \hline
(S) & Electrolysis of 10 L aqueous solution of CuF$_2$ by passing 1 F of electricity. & (4) & pH = 1.0\newline Anode = O$_2$(g)\newline Cathode = Cu(s) \\ \hline
\end{tabular}}\par
\begin{qoptions}
\opt{A}{P $\to$ 1;Q $\to$ 2;R $\to$ 4;S $\to$ 3}
\opt{B}{P $\to$ 2;Q $\to$ 3;R $\to$ 4;S $\to$ 1}
\opt{C}{P $\to$ 1;Q $\to$ 2;R $\to$ 3;S $\to$ 4}
\opt{D}{P $\to$ 4;Q $\to$ 2;R $\to$ 3;S $\to$ 1}
\end{qoptions}
\qn{9} Match the list\par
{\footnotesize\setlength{\tabcolsep}{1.5pt}\noindent\begin{tabular}{|c|L{4.6cm}|c|c|}
\hline
 & \multicolumn{1}{c|}{List-I} & \multicolumn{2}{c|}{List-II} \\ \hline
(P) & Ratio of rate of appearance (gm/min) of SO$_3$(g) and rate of disappearance (gm/min) of SO$_2$(g) for : 2SO$_2$(g) + O$_2$(g) $\longrightarrow$ 2SO$_3$(g) ; would be & (1) & 0 \\ \hline
(Q) & Based on following the observations for a reaction A $\xrightarrow{k}$ product\newline [A]$_0$ (in M) : 0.0677 \ \ 0.136 \ \ 0.272\newline t$_{1/2}$ (in min) : 2 \ \ 4 \ \ 8\newline the order of reaction would be & (2) & 3 \\ \hline
(R) & Time (min) taken by a reaction to decrease the concentration of the reactant(R) from 0.2 mole/litre to 0.025 mole/litre if its rate constant corresponding to rate of disappearance of R is 0.693 min$^{-1}$, would be & (3) & 1.25 \\ \hline
(S) & If emf at 25$^\circ$C of the following Galvanic cell\newline Pt $|$ H$_2$(g) $|$ CH$_3$COOH(aq.) $\|$ HCOOH(aq.) $|$ H$_2$(g) $|$ Pt\newline \hspace*{4.6em}(C$_1$)\hspace*{4.2em}(C$_2$)\newline is 0.0177 volt then ratio of [H$^+$] in anodic compartment and cathodic compartment would be approximately & (4) & 2 \\ \hline
\end{tabular}}\par
\begin{qoptions}
\opt{A}{P $\to$ 3;Q $\to$ 1;R $\to$ 2;S $\to$ 4}
\opt{B}{P $\to$ 2;Q $\to$ 2;R $\to$ 3;S $\to$ 1}
\opt{C}{P $\to$ 1;Q $\to$ 3;R $\to$ 3;S $\to$ 4}
\opt{D}{P $\to$ 3;Q $\to$ 2;R $\to$ 1;S $\to$ 4}
\end{qoptions}
\qn{10} Match the List-I with List-II\par
{\footnotesize\setlength{\tabcolsep}{1.5pt}\noindent\begin{tabular}{|c|L{1.8cm}|c|L{4.8cm}|}
\hline
\multicolumn{2}{|c|}{List-I} & \multicolumn{2}{c|}{List-II} \\ \hline
(P) & Tyrosine & (1) & Gives effervescence with NaHCO$_3$ \\ \hline
(Q) & O$-$Toluidine & (2) & Gives white precipitate with Bromine H$_2$O \\ \hline
(R) & Mesityl oxide & (3) & Gives alkali soluble product on reaction with Benzene sulphonyl chloride \\ \hline
(S) & Buna-N & (4) & Decolourizes cold KMnO$_4$ \\ \hline
 & & (5) & Sodium fusion extract of the compound on boiling with FeSO$_4$, followed by acidification gives prussian blue colour \\ \hline
\end{tabular}}\par
\begin{qoptions}
\opt{A}{P $\to$ 2;Q $\to$ 5;R $\to$ 1;S $\to$ 1,2,3}
\opt{B}{P $\to$ 3;Q $\to$ 1;R $\to$ 5;S $\to$ 1,4}
\opt{C}{P $\to$ 1;Q $\to$ 3;R $\to$ 4;S $\to$ 4,5}
\opt{D}{P $\to$ 4;Q $\to$ 2;R $\to$ 3;S $\to$ 2,3}
\end{qoptions}
\sect{SECTION-II}
\qn{1} 0.1 mole of white phosphorus reacts with 0.6 mole SOCl$_2$ to form a pungent smelling gas(A) alongwith other products. Gas(A) is reacted with 0.2 mole KMnO$_4$ in acidic medium and the product mixture is treated with excess BaCl$_2$ solution. A white precipitate insoluble in dilute acids is formed. Let the weight of precipitate formed be x g. Find the value of $\dfrac x3$.
\qn{2} Consider following precipitates :\par
CuS, ZnS, HgS, Ni(CN)$_2$, AgCN, AgCl, CuCN, Pb(CN)$_2$, FeS, MnS\par
Let,\par
x = number of precipitates soluble in aq.KCN\par
y = number of precipitates soluble in dil.HCl\par
Find $\dfrac{|x-y|}{2}$
\qn{3} Consider following minerals/ores :\par
Malachite, Calamine, Siderite, Kaolinite, Cryolite, Sphalerite, Beryl, Kernite\par
Let the number of types of elements present in all of the above be x. Find $\dfrac x2$
\qn{4} Total number of mixtures among the following that can form a maximum boiling azeotrope is/are -\par
(i) Carbon tetrachloride + Methanol \hfill (ii) Carbon disulphide + Acetone\par
(iii) Benzene + Toluene \hfill (iv) Benzoic acid + Aniline\par
(v) Chloroform + Ethanol \hfill (vi) Water + Ethanol\par
(vii) Bromobenzene + Chlorobenzene \hfill (viii) Bromoethane + Chloroethane\par
(ix) Suphuric acid + Water
\qn{5} In an aqueous solution, mole fraction of solute is 0.2 and its density is d$_1$(g/ml). It is diluted such that mole fraction of solute becomes 0.1 and density becomes d$_2$ g/ml. The density of new solution is 20\% less than initial solution and molarity of new solution\ \ is almost half of the original solution. If molar mass of the solute is `X' gm/mol, then find X.
\qn{6} The molecular weight of the major product of the reaction is \textbf{X} and number of localized lone pair in final product are \textbf{Y}. Then X $\div$ Y is ?\par
\qfig{p1_chem_s3_q06_1}
[Given : Atomic weight of H =1, C = 12, N = 14 and Br = 80 ]
\qn{7} 1.1 mole of (A) was reacted according to below reaction. `C' is the gas produced which turns lime water milky. The amount of `C' formed in gram is \textbf{X}. The value of \textbf{X}/14 is\par
(Given At. Wt. of C, O, H are 12, 16, 1 gm, respectively)\par
\qfig{p1_chem_s3_q07_1}
\qn{8} \textbf{X} is the number of below compounds which are more basic than aniline?\par
o-toluidine, o-anisidine, m-anisidine, p-anisidine, p-chloroaniline, m-toluidine,\par
\qfig{p1_chem_s3_q08_1}
\textbf{Y} is the number of below compounds which are more acidic than benzoic acid\par
\qfig{p1_chem_s3_q08_2}
Picric acid, Hydrofluoric acid, Phosphoric acid (pka$_1$) , Hypochlorous acid\par
The value of (X/Y) is....
\raggedcolumns
\papersec{1}
\partsec{3}{MATHEMATICS}
\sect{SECTION-I(i)}
\qn{1} A letter lock consists of three rings marked with 15 different letters. If N denotes the number of ways which it is possible to make unsuccessful attempts to open the lock, then
\begin{qoptions}
\opt{A}{N is divisible by 482}
\opt{B}{N is product of 3 distinct prime numbers.}
\opt{C}{N is product of 4 distinct prime numbers.}
\opt{D}{The number of positive integer divisors of N is 8}
\end{qoptions}
\qn{2} Twenty children are standing in a line outside a ticket window at\ \ Appu Ghar in New Delhi. Ten of these children have a one rupee coin each and the remaining 10 have a two rupee coin each. The entry ticket is priced at Re. 1. If all the arrange-ments of the 20 children are equally likely, the probability that the 10$^{\mathrm{th}}$ child will be the first to wait for change is (assume that cashier has no change to begin with).
\begin{qoptions}
\opt{A}{0}
\opt{B}{$<\dfrac12$}
\opt{C}{$>$ 1/2}
\opt{D}{1/2}
\end{qoptions}
\qn{3} Let $\vec a=2\hat i-\hat j+\hat k,\vec b=\hat i+2\hat j-\hat k$ and $\vec c=\hat i+\hat j-2\hat k$ be three vectors. A vector in the plane of $\vec b$ and $\vec c$ whose projection on $\vec a$ is of magnitude $\sqrt{2/3}$ is
\begin{qoptions}
\opt{A}{2$\hat i$ + 3$\hat j-$ 3$\hat k$}
\opt{B}{2$\hat i$ + 3$\hat j$ + 3$\hat k$}
\opt{C}{$-$2$\hat i-\hat j$ + 5$\hat k$}
\opt{D}{2$\hat i+\hat j$ + 5$\hat k$}
\end{qoptions}
\qn{4} Let a and c be two complex numbers. Then there is atleast one complex number z such that $|$z $-$ a$|$ + $|$z + a$|$ = 2 $|$c$|$ if
\begin{qoptions}
\opt{A}{$|$c$|$ = $|$a$|$}
\opt{B}{$|$c$|$ $<$ $|$a$|$}
\opt{C}{$|$c$|$ $>$ $|$a$|$}
\opt{D}{2 $|$c$|$ $<$ $|$a$|$}
\end{qoptions}
\qn{5} Let F be the foot of perpendicular from P(1, 2, $-$3) on the line L = $\dfrac{x+1}2=\dfrac{y-3}{-2}=\dfrac z{-1}$. Given Q (x$_1$, y$_1$, z$_1$) and R (x$_2$, y$_2$, z$_2$) are points at a distance of 3 units from F on line L. Let $\pi_1$ and $\pi_2$ be two planes passing through Q and R respectively and perpendicular to the line of intersection of the planes 9x $-$ 7y + 6z + 48 = 0 and x + y $-$ z = 7. Then which of the following statement(s) is(are) correct?
\begin{qoptions}
\opt{A}{The co-ordinates of F is (1, 1, $-$ 1).}
\opt{B}{The value of $\displaystyle\sum_{i=1}^{2}\left(x_i^2+y_i^2+z_i^2\right)$ equals 24.}
\opt{C}{The distance between planes $\pi_1$ and $\pi_2$ equals $\dfrac{88}{\sqrt{482}}$.}
\opt{D}{The normal vector of planes $\pi_1$ and $\pi_2$ is parallel to the vector $\hat i$ + 15$\hat j$ + 16$\hat k$.}
\end{qoptions}
\qn{6} Upon X-Y plane, the path defined by the equation $\dfrac1{x^m}+\dfrac1{y^m}+\dfrac k{(x+y)^n}$ = 0, is
\begin{qoptions}
\opt{A}{a parabola if m = $-\dfrac12$, k = $-$1, n = 0}
\opt{B}{a hyperbola if m = 1, k = $-$1, n = 0}
\opt{C}{a\ \ parabola if m = k = n = 1}
\opt{D}{a pair of lines if m = k = $-$1, n = 1}
\end{qoptions}
\sect{SECTION-I(ii)}
\qn{7} {\footnotesize\setlength{\tabcolsep}{1.5pt}\noindent\begin{tabular}{|c|L{4.2cm}|c|L{1.6cm}|}
\hline
\multicolumn{2}{|c|}{List -I} & \multicolumn{2}{c|}{List -II} \\ \hline
(P) & Locus of centre of circles touching circles x$^2$ + y$^2$ $-$ 4x $-$ 4y = 0 internally and x$^2$+y$^2-$6x$-$6y+17 = 0 externally & (1) & straight line \\ \hline
(Q) & The locus of the point (3h $-$ 2, 3k), where (h, k) lies on the circle x$^2$ + y$^2$ $-$ 2x $-$ 4y $-$ 4 = 0 & (2) & Circle \\ \hline
(R) & Locus of centre of circles touching two circle x$^2$ + y$^2$ + 2x = 0 and x$^2$ + y$^2$ $-$ 6x + 5 = 0 externally is & (3) & Ellipse \\ \hline
(S) & The extremities of a diagonal of a rectangle are (0, 0) and (4, 4). The locus of the extremities of the other diagonal is & (4) & Hyperbola \\ \hline
\end{tabular}}\par
\begin{qoptions}
\opt{A}{P $\to$ 3;Q $\to$ 2;R $\to$ 4;S $\to$ 2}
\opt{B}{P $\to$ 1;Q $\to$ 4;R $\to$ 3;S $\to$ 2}
\opt{C}{P $\to$ 4;Q $\to$ 2;R $\to$ 3;S $\to$ 4}
\opt{D}{P $\to$ 2;Q $\to$ 4;R $\to$ 1;S $\to$ 3}
\end{qoptions}
\qn{8} Let A be a set containing n elements. P and Q are two subsets of A. The number of ways of Choosing P and Q so that\par
{\footnotesize\setlength{\tabcolsep}{1.5pt}\noindent\begin{tabular}{|c|L{3.0cm}|c|L{1.4cm}|}
\hline
\multicolumn{2}{|c|}{List -I} & \multicolumn{2}{c|}{List -II} \\ \hline
(P) & P $\cup$ Q = A & (1) & 2$^n$ \\ \hline
(Q) & P $\cup$ Q = $\phi$ & (2) & 3$^n$ \\ \hline
(R) & P $\cap$ Q is a singleton set & (3) & n.3$^{n-1}$ \\ \hline
(S) & P $\cap$ Q = $\phi$ & (4) & 4$^n-$ 1 \\ \hline
\end{tabular}}\par
\begin{qoptions}
\opt{A}{P $\to$ 3;Q $\to$ 2;R $\to$ 4;S $\to$ 2}
\opt{B}{P $\to$ 1;Q $\to$ 4;R $\to$ 3;S $\to$ 2}
\opt{C}{P $\to$ 4;Q $\to$ 2;R $\to$ 3;S $\to$ 4}
\opt{D}{P $\to$ 2;Q $\to$ 4;R $\to$ 3;S $\to$ 2}
\end{qoptions}
\qn{9} {\footnotesize\setlength{\tabcolsep}{1.5pt}\noindent\begin{tabular}{|c|L{4.4cm}|c|c|}
\hline
\multicolumn{2}{|c|}{List -I} & \multicolumn{2}{c|}{List -II} \\ \hline
(P) & Let f(x) is a derivable function satisfying f(x) = $\displaystyle\int_0^x$ e$^t$ sin(x $-$ t)dt and g(x) = f$''$(x) $-$ f(x), where f$''$(x) = 2$^{\mathrm{nd}}$ order derivative with respect to x. Then the possible integers in the range of g(x) is & (1) & 1 \\ \hline
(Q) & If the substitution x = tan$^{-1}$ (t) transforms the differential equation $\dfrac{d^2y}{dt^2}$ + xy$\dfrac{dy}{dx}$ + sec$^2$x = 0 into a differential equation (1 + t$^2$)$\dfrac{d^2y}{dt^2}$ + $\left(2t+y\tan^{-1}(t)\right)\dfrac{dy}{dt}$ = k then k is equal to & (2) & 0 \\ \hline
(R) & If a$^2$ + b$^2$ = 1 then a$^3$b $-$ ab$^3$ can be equal to & (3) & $-$1 \\ \hline
(S) & system of equation x $-$ $\lambda$y $-$ z = 0 \newline $\lambda$x $-$ y $-$ z = 0 \newline x + y $-$ z = 0 has a unique solution $\lambda$, can be & (4) & 2 \\ \hline
\end{tabular}}\par
\begin{qoptions}
\opt{A}{P $\to$ 1,2;Q $\to$ 1;R $\to$ 2;S $\to$ 2,4}
\opt{B}{P $\to$ 1,2;Q $\to$ 2,4;R $\to$ 1;S $\to$ 1,2}
\opt{C}{P $\to$ 2,4;Q $\to$ 2;R $\to$ 3,4;S $\to$ 1,2}
\opt{D}{P $\to$ 1;Q $\to$ 2;R $\to$ 1,2;S $\to$ 2,4}
\end{qoptions}
\qn{10} {\footnotesize\setlength{\tabcolsep}{1.5pt}\noindent\begin{tabular}{|c|L{4.4cm}|c|c|}
\hline
\multicolumn{2}{|c|}{List -I} & \multicolumn{2}{c|}{List -II} \\ \hline
(P) & If the period of sin$^2$2x + cos$^4$2x + 2 is T; then $\dfrac\pi T$ is equal to & (1) & $-$2 \\ \hline
(Q) & If 4$^x-$ 2$^{x+2}$ + 5 + $||\lambda-1|-3|$ = $|\sin y|$ , x, y, $\lambda\in$ R then the possible values of $\lambda$ is/are & (2) & 2 \\ \hline
(R) & Let f(x) = min $\{x^2, 2\}$, then the possible integral values of `a' for which f(x) $<$ a for atleast one real `x' is/are & (3) & 4 \\ \hline
(S) & Let f(x) = $\begin{cases}|2x-1|,&-2\leqslant x\leqslant1\\x^2-4,&1<x\leqslant8\end{cases}$, then possible integers in the domain of y = f(f(x)) is & (4) & 6 \\ \hline
\end{tabular}}\par
\begin{qoptions}
\opt{A}{P $\to$ 3;Q $\to$ 1,3;R $\to$ 1,2,3;S $\to$ 2,3,4}
\opt{B}{P $\to$ 3;Q $\to$ 1,3;R $\to$ 2,3,4;S $\to$ 1,2,3}
\opt{C}{P $\to$ 1,3;Q $\to$ 1;R $\to$ 3,4;S $\to$ 1,2,3}
\opt{D}{P $\to$ 2,3,4;Q $\to$ 1,2,3;R $\to$ 3;S $\to$ 1,3}
\end{qoptions}
\sect{SECTION-II}
\qn{1} If $\displaystyle\int e^{\tan\theta}(1-\tan\theta)^2\cos^2\theta\,d\theta=Ae^{\tan\theta}\cos^B\theta+K$, then the value of A + B is
\qn{2} The value of $\displaystyle\lim_{n\to\infty}\sum_{r=1}^{n}\left(\sum_{k=0}^{r-1}\dfrac1{5^{n-1}}\,{}^nC_r\,{}^rC_k3^k\right)$ is equal to
\qn{3} If z$_1$, z$_2$ are complex number such that z$_1^3-$3z$_1$z$_2^2$ = 2 and 3z$_1^2$z$_2-$z$_2^3$ = 11 then $|z_1^2+z_2^2|$ =
\qn{4} Let S be the focus of the parabola y$^2$ = 8x and let PQ be the common chord of the circle x$^2$ + y$^2$ $-$ 2x $-$ 4y = 0 and the given parabola. The area of the triangle PQS is
\qn{5} If the line $\dfrac{x-1}{15}=\dfrac{y+1}{16}=\dfrac{z-1}n$ intersect the curve 6x$^2$ + 5y$^2$ = 1, z = 0 ; then 10n$^2-$ 20n + k = 0, where the value of K $-$ 2630 is
\qn{6} Let A be a 2 $\times$ 3 matrix whereas B be a 3 $\times$ 2 matrix. If det (AB) = 4, then the value of det (BA) is
\qn{7} If y$^2$ = P(x) is a polynomial of degree 3, then 2$\dfrac d{dx}\left(y^3\dfrac{d^2y}{dx^2}\right)$ equals $\dfrac k5$P(x)$\dfrac{d^3P(x)}{dx^3}$ then k is equal to
\qn{8} A sequence is obtained by deleting all perfect squares from set of natural numbers. The remainder when the 2003$^{\mathrm{rd}}$ term of new sequence is divided by 2048 , is
\raggedcolumns
\papersec{2}
\partsec{1}{PHYSICS}
\sect{SECTION-I(i)}
\qn{1} On the superposition of two harmonic oscillations represented by x$_1$ = a sin ($\omega$t + $\phi_1$) and x$_2$ = a sin ($\omega$t + $\phi_2$) a resulting oscillation with the same time period and amplitude is obtained. The value of $\phi_1-\phi_2$ is :-
\begin{qoptions}
\opt{A}{120$^\circ$}
\opt{B}{90$^\circ$}
\opt{C}{60$^\circ$}
\opt{D}{15$^\circ$}
\end{qoptions}
\qn{2} An uncharged sphere of metal is placed inside a charged parallel plate capacitor. The electric lines of force will look like
\begin{qoptions}
\opt{A}{\qfigi{p2_phy_s1_q02_A}}
\opt{B}{\qfigi{p2_phy_s1_q02_B}}
\opt{C}{\qfigi{p2_phy_s1_q02_C}}
\opt{D}{\qfigi{p2_phy_s1_q02_D}}
\end{qoptions}
\qn{3} An electromagnetic wave travelling along z-axis is given as: $\vec E=\vec E_0\cos(kz-\omega t)$. Choose the \textbf{INCORRECT} options from the following;
\begin{qoptions}
\opt{A}{The associated magnetic field is given as $\vec B=\dfrac1c\hat k\times\vec E=\dfrac1\omega\left(\vec k\times\vec E\right)$}
\opt{B}{The electromagnetic field can be written in terms of the associated magnetic field as $\vec E=c\left(\vec B\times\hat k\right)$}
\opt{C}{$\hat k\cdot\vec E=0,\hat k\cdot\vec B=0$}
\opt{D}{$\hat k\times\vec E=0,\hat k\times\vec B=0$}
\end{qoptions}
\qfullbegin
\qn{4} An $\alpha$-particle having a de Broglie wavelength $\lambda_i$ collides with a stationary carbon nucleus. The $\alpha$-particle moves off in a different direction as shown below.\par
\qfig{p2_phy_s1_q04_1}
\qfullsplit
After the collision, the de Broglie wavelength of the $\alpha$-particle and the carbon nucleus are $\lambda_f$ and $\lambda_e$ respectively. Which of the following relations about de Broglie wavelength is correct ?
\begin{qoptions}
\opt{A}{$\lambda_i<\lambda_f$}
\opt{B}{$\lambda_i>\lambda_f$}
\opt{C}{$\lambda_f=\lambda_e$}
\opt{D}{$\lambda_i=\lambda_e$}
\end{qoptions}
\qfullend
\sect{SECTION-I(ii)}
\qn{5} Consider the following constraint motion. If all the strings remain tight then choose the correct statement(s)\par
\qfig{p2_phy_s2_q05_1}
\begin{qoptions}
\opt{A}{Velocity of ring (1) is $\dfrac{25}4$m/s towards left.}
\opt{B}{Velocity of ring (1) is $\dfrac{25\sqrt3}8$m/s towards left.}
\opt{C}{Velocity of ring (2) is $\dfrac{25}4$m/s towards left.}
\opt{D}{Velocity of ring (2) is 4 m/s towards left.}
\end{qoptions}
\qn{6} S is a solid neutral conducting sphere. A point charge q of 1 $\times$ 10$^{-6}$C is placed at point A. C is the centre of sphere and AB is a tangent. BC = 3m and AB = 4m.\par
\qfig{p2_phy_s2_q06_1}
\begin{qoptions}
\opt{A}{The electric potential of the conductor is 1.8 kV.}
\opt{B}{The electric potential of the conductor is 2.25 kV.}
\opt{C}{The electric potential at B due to induced charges on the sphere is $-$ 0.45 kV.}
\opt{D}{The electric potential at B due to induced charges on the sphere is 0.45 kV.}
\end{qoptions}
\qn{7} Two concentric conducting spherical shells of radii R and 2R having charges q$_\mathrm A$ and q$_\mathrm B$ and potential 2v and 3v/2 respectively. Now shell B is earthen by closing switch S. Let charges on spherical shells A \& B become q$'_\mathrm A$ \& q$'_\mathrm B$ respectively, then\par
\qfig{p2_phy_s2_q07_1}
\begin{qoptions}
\opt{A}{q$_\mathrm A$/q$_\mathrm B$ = $\frac12$ (magnitude)}
\opt{B}{q$'_\mathrm A$/q$'_\mathrm B$ = 1 (magnitude)}
\opt{C}{potential of A after earthing becomes $\dfrac{3v}2$}
\opt{D}{potential difference between A and B after earthing becomes $\dfrac v2$}
\end{qoptions}
\qn{8} If a graph of Binding energy (B/A) per nucleon versus mass number (A) looks like as shown in figure. Then using the curve choose \textbf{CORRECT} option(s).\par
\qfig{p2_phy_s2_q08_1}
\begin{qoptions}
\opt{A}{Fusion of two nuclei with mass numbers 30 and 45 will have Q value 99 MeV}
\opt{B}{Fission of nuclei with mass number 80 into two nuclei of equal mass number will have Q value 256 MeV}
\opt{C}{Fission of nuclei with mass number 150 into two nuclei of equal mass number will release energy}
\opt{D}{Fission of nuclei with mass number 80 into two nuclei of equal mass number will release energy}
\end{qoptions}
\qn{9} An inductor of inductance L = 10H carrying initial current 5A is connected to a resistor of resistance 20$\Omega$ and mass 2kg and specific heat 0.5Jkg$^{-1}$K$^{-1}$. Assume that all energy given to resistor is used to increase its temperature .
\begin{qoptions}
\opt{A}{The total change in temperature of resistor will be 125$^\circ$C.}
\opt{B}{The total change in temperature of resistor will be 250$^\circ$C.}
\opt{C}{The total charge passed through the resistor is 2.5 C.}
\opt{D}{The total charge passed through the resistor is infinite.}
\end{qoptions}
\qn{10} A ball of density $\rho$ is dropped onto a horizontal solid surface.\ \ It bounces elastically from the surface and returns to its original position in time t$_1$.\ \ Next, the ball is released from the same height, but this time it strikes the surface of the liquid of density $\rho_L$ ($>\rho$ but less than 2$\rho$), and takes t$_2$ second to come back to its original height:
\begin{qoptions}
\opt{A}{t$_2>$ t$_1$}
\opt{B}{t$_2<$ t$_1$}
\opt{C}{The motion of the ball is not simple harmonic}
\opt{D}{If $\rho=\rho_L$, then the speed of the ball inside the liquid will be independent of its depth. (neglect viscous\ \ force)}
\end{qoptions}
\sect{SECTION-II}
\qn{1} A rigid triangular frame consists of three massless rods of length $\ell$=2.5 m each and point masses of mass m each at vertex B and C respectively. Frame is hanging vertically from point A about which it can rotate freely about an axis xx$'$ which is perpendicular to plane of frame as shown in figure. Point of suspension of frame, i.e. A, is accelerating with constant acceleration a = 3g/4 in horizontal direction and initially frame is at rest w.r.t. support A. Minimum initial angular velocity $\omega$ = x(3)$^{\frac14}$ (in rad/s) provided to system, so that it can complete vertical circular motion in the frame of support A. Calculate the value of x. (take g = 10 m/s$^2$)\par
\qfig{p2_phy_s3_q01_1}
\qn{2} For an equilateral prism, it is observed that when a ray strikes grazingly at one face it emerges grazingly at the other. Find the refractive index of the prism.
\qn{3} The following figure shows a circuit diagram of an ammeter of three different ranges. The specification of the ammeter is\par
\qfig{p2_phy_s3_q03_1}
The resistance of the galvanometer coil is 99 $\Omega$. The range of the galvanometer is 10 mA. Find the $\left(\dfrac{R_3}{R_1+R_2}\right)$.
\qn{4} A monochromatic point source of light is placed at a distance `d' from a small flat metal surface. Photo-electrons are ejected at rate n. If the source is brought nearer to distance d/2, the rate becomes p times of its initial value. Then find the value of p.
\qn{5} A tuning fork P of unknown frequency gives 7 beats in 2 seconds with another tuning fork Q. When Q is moved towards a wall with a speed of 5 m/s, it gives 5 beats per second for a stationary observer located left to it.\ \ On filing, P gives 6 beats per second with Q. What is the frequency (in Hz) of P? Assume speed of sound = 332 m/s.
\qn{6} The friction coefficient between the horizontal surface and each of the block shown in the figure is 0.4. The collision between the blocks is perfectly elastic. If the separation (in cm) between them when they come to rest is x then find the value 8x.\par
\qfig{p2_phy_s3_q06_1}
\qn{7} A uniform rod of length 0.3m, hanging vertically. The free end acquires velocity of v m/s, when it is hit at its free end. Find the maximum value of v in m/s for which the rod does not reach horizontal position.
\qn{8} Two steel wires of radii r and 2r are connected together end to end and tied to a wall as shown.\par
\qfig{p2_phy_s3_q08_1}
The force stretches the combination by $\dfrac{27}4$mm. How far does the junction point A move. (in mm)
\raggedcolumns
\papersec{2}
\partsec{2}{CHEMISTRY}
\sect{SECTION-I(i)}
\qn{1} Which is NOT a redox reaction ? (ignore balancing)
\begin{qoptions}
\opt{A}{Pb$_3$O$_4$ + HNO$_3$ $\to$}
\opt{B}{CaOCl$_2$ + dil. HCl $\to$}
\opt{C}{XeO$_6^{4-}$ + F$^{-1}$ + H$^+$ $\to$}
\opt{D}{Na$_2$S + Na$_2$[Fe(CN)$_5$NO] $\to$}
\end{qoptions}
\qn{2} Consider following structure of N$_2$O$_3$ :\par
\qfig{p2_chem_s1_q02_1}
Choose incorrect option :\par
I. y $>$ z = t $>$ x\par
II. $\alpha$ $>$ $\beta$ $>$ $\gamma$
\begin{qoptions}
\opt{A}{I only}
\opt{B}{II only}
\opt{C}{Both}
\opt{D}{None}
\end{qoptions}
\qn{3} At low pressure, if RT = 2$\sqrt{P\times a}$, then the molar volume occupied by a real gas is, where `a' is vanderwaal's coefficient.
\begin{qoptions}
\opt{A}{$\dfrac{2RT}P$}
\opt{B}{$\dfrac{2P}{RT}$}
\opt{C}{$\dfrac{RT}{2P}$}
\opt{D}{$\dfrac{3RT}P$}
\end{qoptions}
\qfullbegin
\qn{4} The rate\ \ constants for the bromination of several disubstituted stilbenes are given. Find the incorrect relation between rate constants for different X and Y?\par
The rate the table\par
\qfig{p2_chem_s1_q04_1}
\qfullsplit
{\footnotesize\setlength{\tabcolsep}{4pt}\noindent\begin{tabular}{|c|c|c|}
\hline
X & Y & K($\times$10$^3$l/mol-min) \\ \hline
OMe & Me & K$_1$ \\ \hline
OMe & Cl & K$_2$ \\ \hline
OMe & OMe & K$_3$ \\ \hline
OMe & NO$_2$ & K$_4$ \\ \hline
\end{tabular}}\par
\begin{qoptions}
\opt{A}{K$_3$ $>$ K$_1$ $>$ K$_2$}
\opt{B}{K$_1$ $>$ K$_2$ $>$ K$_4$}
\opt{C}{K$_3$ $>$ K$_2$ $>$ K$_4$}
\opt{D}{K$_1$ $>$ K$_3$ $>$ K$_4$}
\end{qoptions}
\qfullend
\sect{SECTION-I(ii)}
\qn{5} Choose correct reaction(s)/order(s)/statement(s)\par
I. 2CrCl$_2$ + 2HCl $\to$ 2CrCl$_3$ + H$_2$\par
II. Cu$_2$CO$_3$ + 2HCl $\to$ CuCl$_2$ + Cu + H$_2$O + CO$_2$\par
III. Cr$^{2+}$ $>$ Fe$^{2+}$ (reducing power)\par
IV. Mn$^{3+}$ $>$ Co$^{3+}$ (oxidising power)\par
V. Na$_2$Cr$_2$O$_7$ is used as a primary standard in volumetric analysis as it is deliquescent\par
VI. V$_2$O$_5$ is amphoteric oxide as it reacts with both acids and alkalies
\begin{qoptions}
\opt{A}{I, III}
\opt{B}{IV, V}
\opt{C}{IV, VI}
\opt{D}{I, II, V, VI}
\end{qoptions}
\qn{6} Consider following complex :\par
[Mn(gly)$_2$ (NO$_2$)(CN)]$^{-1}$\par
Choose incorrect statements regarding this complex :
\begin{qoptions}
\opt{A}{There are 4 geometrical isomers in which NO$_2$ and CN are in cis position w.r.t. each other}
\opt{B}{Splitting energy is less than pairing energy}
\opt{C}{CFSE = $-$0.6$\Delta_0$}
\opt{D}{Mn $-$ C bond will have partial double bond character}
\end{qoptions}
\qn{7} 100 ml of 1 M HCN is titrated with 5 M KMnO$_4$ in acidic medium, A gas (B) is liberated which turns lime water milky, along with an anion which gives brown ring test (central atom in the highest\ \ oxidation state). When this gas is passed through lime water, 10 gms of a white ppt. is obtained. Identify the correct statement(s) :-
\begin{qoptions}
\opt{A}{Normality of the HCN solution would be 10N}
\opt{B}{Volume of the KMnO$_4$ solution at equivalence point would be 40 ml}
\opt{C}{Equivalent mass of HCN in this process is greater than its molecular mass}
\opt{D}{White precipitate formed in this reaction would be CaC$_2$}
\end{qoptions}
\qn{8} Choose the correct statement(s) among the following regarding hydrogen atom?\par
(All graphs given below are indicative in nature, not to scale)
\begin{qoptions}
\opt{A}{For an atomic orbital, radial wave function R(r) depends on principle, azimuthal, and magnetic quantum number.}
\opt{B}{Plot of R$_{(2,0)}$(r) versus r (distance from nucleus), for 2s electron is \qfigi{p2_chem_s2_q08_1}}
\opt{C}{Plot of R$^2_{(2,0)}$(r) versus r (distance from nucleus), for 2s electron is \qfigi{p2_chem_s2_q08_2}}
\opt{D}{Plot of 4$\pi$r$^2$R$^2_{(2,0)}$(r) against r (distance from nucleus) is \qfigi{p2_chem_s2_q08_3}}
\end{qoptions}
\qn{9} Choose the \textbf{CORRECT} statement(s) for given reaction scheme\par
\qfig{p2_chem_s2_q09_1}
In the above reaction scheme, if [Q] is urea and [R] is veronal, then correct statement(s) is/are :-
\begin{qoptions}
\opt{A}{[R] is a derivative of barbituric acid}
\opt{B}{[R] is used as a hypnotic.}
\opt{C}{[R] contains two delocalized lone pair.}
\opt{D}{There will be three C=O groups in [R]}
\end{qoptions}
\qfullbegin
\qn{10} Vanillin (L) C$_8$H$_8$O$_3$ is isolated from vanilla beans. It gives an intense colour with FeCl$_3$ and a positive Tollen's test. It is not steam-distilled and does not react with HCl. It goes through the following steps :\par
\qfig{p2_chem_s2_q10_1}
\qfullsplit
Among the following pair(s) which can be distinguished by aq. NaHCO$_3$ is/are :-
\begin{qoptions}
\opt{A}{M, N}
\opt{B}{N, P}
\opt{C}{L,O}
\opt{D}{P, O}
\end{qoptions}
\qfullend
\sect{SECTION-II}
\qn{1} Consider following compounds :\par
MnO$_2$, HgI$_2$, Cu(CN)$_2$, Cu$_2$[Fe(CN)$_6$], Fe(CN)$_3$, HgSO$_4$.2HgO, Cu$_2$O, MnS, [Ni(CN)$_4$]$^{2-}$\par
Let,\par
x = number of yellow coloured compounds\par
y = number of brown coloured compounds\par
Find $\dfrac{xy}2$
\qn{2} Consider following molecules/ions :\par
SF$_4$, ClF$_3$, N$_2$O, XeO$_6^{4-}$, Al$_2$Cl$_6$, BrF$_5$, C$_6$H$_6$\par
Let\par
x = number of molecules/ion having all bond length equal\par
y = number of polar as well as planar species. Find $\dfrac{y-x}3$
\qn{3} Solid AgNO$_3$ is added to an aqueous solution containing 0.2 M Br$^\ominus$ and 0.1 M I$^\ominus$. Assuming that AgNO$_3$ addition does not change volume of this solution, calculate \% [I$^-$] precipitated when [Br$^-$] just starts precipitating :\ \ K$_\mathrm{sp}$(AgBr) = 10$^{-13}$ M$^2$; k$_\mathrm{sp}$(AgI) = 10$^{-17}$ M$^2$
\qn{4} An ideal gas having initial pressure P, volume V and temperature T is allowed to expand adiabatically and reversibly until its volume becomes 5.66V while its temperature falls to half of initial temperature. The work done by the gas is represented by x times PV. Find the value of x.\par
(10$^{0.3}$ = 2, 10$^{0.48}$ = 3, 10$^{0.75}$ = 5.66)
\qn{5} X constitutes an FCC lattice with Y occupying all octahedral voids and Z occupying all tetrahedral voids. One such isolated unit cell is divided into 8 equal parts, then select total number of correct statements among the following\par
(i)\ \ \ Every equal part contains all X, Y, and Z.\par
(ii)\ \ In smaller part, 4 corners are occupied by X and 4 corners are occupied by Y.\par
(iii) In smaller part, 4 corners are occupied by X, 2 by Y and 2 by Z.\par
(iv) In smaller part, effective numbers are X = $\dfrac12$, Y = $\dfrac12$, Z = 1\par
(v)\ \ Co-ordination number of X with respect of Z in the lattice is 8\par
(vi) Co-ordination number of X with respect of Y in the lattice is 6\par
(vii) Number of next nearest neighbour of Y is 8\par
(viii) If edge length of $\dfrac18$th equal part of the isolated unit cell is b, then the minimum distance between two Y atoms would be $\sqrt2$b
\qn{6} An alkene C$_n$H$_{2n}$ is converted by a reduction procedure to alkane C$_n$H$_{2n+2}$. The \% change in the molecular weight = 3.57. What is the value of n. (rounded off to the nearest integer.)
\qn{7} (a)\ \ Leucine , Lysine , Serine, Cysteine , Proline, Threonine , Arginine , Histidine\par
\hspace*{1.2em}Number of Essential amino acids in above mentioned are X.\par
(b)\ \ In phosphodiester linkage of a Dinucleotide, carbon number (Y) of one sugar and carbon number (Z) of other sugar are joined by phosphate. Then X $\div$ (Y + Z)\ \ is :
\qn{8} The amount of product produced in the following reactions\ \ is \underline{\hspace{1cm}}kg.\par
C$_2$H$_4$(2kg) + H$_2$(2kg) $\xrightarrow{\text{Wilkinson}'\text{s Catalyst}}$ product (90 \% catalytic conversion)
\raggedcolumns
\papersec{2}
\partsec{3}{MATHEMATICS}
\sect{SECTION-I(i)}
\qn{1} $\displaystyle\lim_{x\to0}\dfrac{\ln(\sec(ex)\ \sec(e^2x)...\sec(e^{50}x))}{e^2-e^{2\cos x}}$ is equal to
\begin{qoptions}
\opt{A}{$\dfrac{e^{100}-1}{2(e^2-1)}$}
\opt{B}{$\dfrac{e^{100}-1}{(e^2-1)}$}
\opt{C}{$\dfrac{e^{100}-1}{2e(e^2-1)}$}
\opt{D}{$\dfrac{e^{50}-1}{2e(e^2-1)}$}
\end{qoptions}
\qn{2} Let a be a fixed real number. Consider the equation,\par
(x + 2)$^2$(x + 7)$^2$ + a = 0, x $\in$ R\par
This equation has exactly one double root. Possible value of a is
\begin{qoptions}
\opt{A}{0}
\opt{B}{$\dfrac{-9}2$}
\opt{C}{$-\left(\dfrac52\right)^4$}
\opt{D}{$-$2}
\end{qoptions}
\qn{3} If the equation x$^3-$ 3x + 1 = 0 has three real roots x$_1$, x$_2$, x$_3$ where x$_1$ $<$ x$_2$ $<$ x$_3$ then the value of ($\{$x$_1\}$ + $\{$x$_2\}$ + $\{$x$_3\}$) is, (where $\{\cdot\}$ is fractional part function )
\begin{qoptions}
\opt{A}{$\dfrac32$}
\opt{B}{1}
\opt{C}{2}
\opt{D}{$\dfrac52$}
\end{qoptions}
\qn{4} Equation of a plane passing through the line of intersection of the planes\par
x + 2y + z = 10 and 3x + y $-$ z = 5 and passing through origin is
\begin{qoptions}
\opt{A}{5x + 3z = 0}
\opt{B}{5x $-$ 3z = 0}
\opt{C}{5x + 4y + 3z = 0}
\opt{D}{5x $-$ 4y + 3z = 0}
\end{qoptions}
\sect{SECTION-I(ii)}
\qn{5} Given, f(x) $-$ $\mu\displaystyle\int_0^{\pi/4}$ sin$^2$xcos$^{-5}$x\ \ cos t f(t)dt = sin$^2$xcos$^{-5}$x (where $\mu$ is real constant) then
\begin{qoptions}
\opt{A}{f(x) = $\dfrac3{3-\mu}$sin$^2$xcos$^{-5}$x}
\opt{B}{f(x) = $\dfrac3\mu$sin x cos x}
\opt{C}{f(x) is periodic with period $\pi$}
\opt{D}{f(x) is periodic with period 2$\pi$}
\end{qoptions}
\qn{6} A differentiable function f : R $\longrightarrow$ R satisfies following relations\par
(i) f(x) $-$ f(y) = (x $-$ y)f$'\left(\dfrac{x+y}2\right)$, $\forall$ x, y $\in$ R\par
(ii) f(0) = f(1) = 0, f(2) $>$ 0\par
Which of the following is/are true
\begin{qoptions}
\opt{A}{f(x) is an onto function}
\opt{B}{f(x) has an extrema at x = $\dfrac12$}
\opt{C}{f(x) is a many one function}
\opt{D}{$\displaystyle\lim_{x\to-\infty}$ f(x) = $-\infty$}
\end{qoptions}
\qn{7} Let $\hat a$ and $\hat c$ be unit vectors and $\left|\vec b\right|$ = 4. If the angle between $\hat a$ and $\hat c$ is cos$^{-1}\left(\dfrac14\right)$ and $\vec b-2\hat c=\lambda\hat a$ then the value of $\lambda$ can be :
\begin{qoptions}
\opt{A}{2}
\opt{B}{$-$3}
\opt{C}{3}
\opt{D}{$-$4}
\end{qoptions}
\qn{8} Which of the following is/are true ?
\begin{qoptions}
\opt{A}{If z is a complex number satisfying $|$z $-$ 2 $-$ i$|$ = $|$z$|\left|\sin\left(\dfrac\pi4-\arg z\right)\right|$ then locus of z is a parabola}
\opt{B}{If z is a complex number satisfying $\left|z-\cos^{-1}(\cos)2\right|-\left|z-\sin^{-1}(\sin12)\right|$ = 8($\pi-$3) then locus of z is a hyperbola}
\opt{C}{If z is a complex number satisfying z$^2$ + az + a$^2$ = 0 then $|$arg z$|$ can be $\dfrac{2\pi}3$ (a $>$ 0)}
\opt{D}{If z is a complex number satisfying z$^2$ + az + a$^2$ = 0 then arg(z) = $\dfrac\pi3$ (a $\in$ R)}
\end{qoptions}
\qn{9} The differential equation (xy$^3$ (1 + cosx) $-$ y) dx + xdy = 0 represents the curve\par
$\dfrac{x^2}{2y^2}=\dfrac{x^3}b$ + x$^2$ sin x + cx cos x $-$ d sin x + k then, (where b, c, d $\in$ Z (Z is set of integers \& k is constant of integration))
\begin{qoptions}
\opt{A}{d = 2}
\opt{B}{c = 3}
\opt{C}{b = 3}
\opt{D}{c = 2}
\end{qoptions}
\qn{10} If $\displaystyle\int_0^\pi(\sin^3x)e^{-\sin^2x}dx=\alpha-\dfrac\beta e\int_0^1\sqrt t\,e^t\,dt$ then
\begin{qoptions}
\opt{A}{$\alpha$ = 2}
\opt{B}{$\beta$ = $-$3}
\opt{C}{$\alpha$ + $\beta$ = 5}
\opt{D}{$\alpha$ $-$ $\beta$ = 5}
\end{qoptions}
\sect{SECTION-II}
\qn{1} If f(x) = $\left([\{x\}]\tan^{-1}\left(\dfrac{x^2-3x-1}{x^2-3x+5}\right)+3-x^7\right)^{\frac17}$\par
Where [k] and $\{$k$\}$ denotes greatest integer and functional part of k respectively, then the value of f$^{-1}$(50) $-$ f(50) + f(f(100)) , is
\qn{2} The number of 5 digit numbers that can be formed using the digits 1, 2, 3, 4, 5, 6, 7, 8, 9 in which one digit appears once and two digits appear twice (eg., 44533) is k, then the value of $\dfrac k{1890}$ is equal to
\qn{3} If sum of the values of x satisfying the equation\par
$\sqrt2\sin\left(x+\dfrac\pi4\right)\left(\sec^2x-(1+\sqrt3)\tan x+2+\sqrt3\right)=\sqrt{9+18\sin x\cos x}$\par
in [0, 2$\pi$] is $\left(\dfrac pq\right)\pi$, where p and q are co-prime, then $|$p $-$ q$|$ is equal to
\qn{4} If the least positive value of x satisfying the equation $\dfrac{\sin x}{\cos3x}+\dfrac{\sin3x}{\cos9x}+\dfrac{\sin9x}{\cos27x}$ = 0 is $\dfrac\pi k$ where k $\in$ N then $\left(\dfrac k5\right)$ is equal to
\qn{5} If $\dfrac{\cos x+\cos y+\cos z}{\cos(x+y+z)}$ = 2 and $\dfrac{\sin x+\sin y+\sin z}{\sin(x+y+z)}$ = 2\par
then the value of cos(x + y) + cos(y + z) + cos(z + x) is equal to :\par
(where x, y, z $\in$ R)
\qn{6} Let (2x$^2$ + 3x + 4)$^{10}$ = $\displaystyle\sum_{r=0}^{20}$ a$_\mathrm r$x$^\mathrm r$ where a$_0$, a$_1$, a$_2$, ....., a$_{20}$ are constant then the value of $\dfrac{a_7}{a_{13}}$ is equal to
\qn{7} The greatest possible integral value of $\dfrac{b-a}{\tan^{-1}b-\tan^{-1}a}$ , (where 0 $<$ a $<$ b $<$ $\sqrt3$) is
\qn{8} If $\displaystyle\int_0^{\pi/2}\dfrac{dx}{\left(\sqrt{\cos x}+\sqrt{\sin x}\right)^4}$ = J then $\dfrac1J$ is
\end{multicols}
\end{document}
